\section*{De la clôture généalogique à la réalisation physique}

La thèse de cet ouvrage acquiert quatre clôtures coordonnées et irréductibles :
généalogique, opératorielle, matricielle et physique. Il ne s'agit pas de juxtaposer une
théorie des ensembles, une théorie des algèbres d'opérateurs ou un formalisme
quantique à une construction numérique préalable. Les quatre développements naissent de
la même préséance causale : l'APP construit la double réalisation positionnelle et
arithmétique ; le TRIT orienté distingue retour, présent torsionnel et ouverture ;
le TPK conserve valeur, feuillet, orientation, résidu, quotient, retenue, bord,
route, registre, mémoire, survie, phase et jauge. Ce n'est qu'ensuite que se forment
la section compatible, ses lecteurs et ses représentations. Le diagramme directeur
est donc
\[
\begin{gathered}
 \mathrm{APP}
 \longrightarrow \mathrm{TRIT}_{\rm orientado}
 \longrightarrow \mathrm{TPK}
 \longrightarrow \mathcal X_\infty^{\rm cal},\\[2mm]
 \mathcal X_\infty^{\rm cal}
 \xrightarrow{\;\mathfrak R\;}\mathbb R,
 \qquad
 \mathcal X_\infty^{\rm cal}\times\mathcal C_{\rm exc}
 \xrightarrow{\;\mathfrak E\;}\mathcal E_{\rm exc},\\[2mm]
 C(X_m)\xrightarrow{\;\operatorname{diag}\;}M_{9^m}(\mathbb C)
 \xrightarrow{\;a\mapsto a\otimes I_9\;}
 \mathrm{UHF}(9^\infty)
 \xrightarrow{\;\pi_\tau(\,\cdot\,)''\;}R_{\rm hyp},\\[2mm]
 \mathsf P_{\rm CT}^{+}
 \xrightarrow{\;\mathfrak R_{\rm MD}^{\rm disc,+}\;}
 \mathbf{SpecRoute}\times\mathsf{Hilb}_{\mathbb R,J}
 \times\mathbf B\mathfrak M_D^+\times\mathbf B G_C.
\end{gathered}
\]
Les quatre lignes ne forment pas une unique chaîne linéaire. L'évaluation réelle et la
projection exceptionnelle sont des lectrices sœurs de la section enrichie ; la
branche traciale passe par les algèbres cylindriques et leurs inclusions unitaires ; la
réalisation physique part de routes préparées munies des structures
exigées par chaque composante. Cette séparation des domaines empêche de déduire le facteur de
von Neumann directement d'une section, d'identifier la projection
exceptionnelle à une action sur les chiffres ou de présenter une grandeur physique comme
s'il s'agissait d'un état TPK sans morphisme de réalisation.

\subsection*{Clôture généalogique : extensionnalité, infini et décision}

La première clôture détermine ce que signifie conserver un objet à travers
toutes les profondeurs. Étant donné un système inverse d'espaces finis
\((X_m,p_m)\), une section \(x=(x_m)_m\) n'est pas caractérisée uniquement par
l'égalité de ses valeurs publiées. Le théorème de descente par fibres fixe
les conditions dans lesquelles un lecteur fini descend à un quotient et
sépare surjectivité, constance sur les fibres et unicité de l'application descendue.
Le lemme de König garantit l'existence d'une branche lorsque les niveaux sont
finis, non vides et finement ramifiés ; l'unicité exige, en outre, le
caractère complet qui sélectionne un unique prolongement compatible.

Sur une section fixée, les cylindres décroissants fournissent une
réalisation exacte du passage \(0\!-1\!-\!\infty\) : le support se contracte, la
hauteur augmente, la masse demeure unitaire et les mesures cylindriques convergent
faiblement-* vers \(\delta_x\). Ce résultat coordonne le delta de Kronecker fini,
les espérances conditionnelles entre niveaux et le delta de Dirac comme mesure
ponctuelle, sans identifier ce dernier à un opérateur différentiel. La
fondation des préfixes empêche les chaînes descendantes infinies dans chaque
présentation finie ; la coinduction contrôle le prolongement compatible. Les
ordinaux de von Neumann sont relevés avec une fibre de mémoire de route, et la
cardinalité apparaît comme une décatégorification qui ne retient pas toute la généalogie.

La même distinction oblige à séparer puissance cylindrique, puissance fermée et
puissance pleine, ainsi que choix ensembliste, choix naturel, choix
algébrique et choix effectif. Le codage dans ZF/ZFC conserve les
constructions et permet de les confronter à l'extensionnalité, à Fondation,
à Puissance et à Choix ; il ne fait pas de la droite réelle ni de l'univers ensembliste
l'origine d'APP--TRIT--TPK\@. Dans le domaine effectif, la décidabilité de chaque
fenêtre finie, l'existence de la limite, l'unicité d'une instance et
l'impossibilité d'un décideur total uniforme sont quatre assertions distinctes.
La réduction de Godel--Turing n'affecte que la quatrième : elle n'invalide pas les sélecteurs
particuliers déjà construits. De même, la confrontation Godel--Cohen sépare
l'extension \(M\mapsto M[G]\) de l'enrichissement interne d'un état et
évite de confondre le franchissement des denses de profondeur avec la généricité sur un
modèle. Les cylindres boréliens, les hiérarchies de comparaison et le passage
projectif de minimax sont introduits avec leurs hypothèses de compacité,
de continuité et de cohérence, non comme des conséquences automatiques de la notation
inverse.

Cette clôture possède également une formulation informationnelle. Si une mise à jour
\(U\) de l'état complet est bijective, alors
\(H(X\mid U(X))=0\), et le terme logique idéal de Landauer s'annule. Si un
lecteur \(q\) ne conserve que l'ombre visible, l'information écartée est
localisée dans \(H(X\mid q(X))\). Dans la formulation opératorielle, un automorphisme
préservant la trace conserve l'entropie, tandis qu'une espérance conditionnelle
préservant la trace \(E_N:M\to N\) satisfait
\[
 S_\tau(E_Nh)-S_\tau(h)=D_\tau(h\Vert E_Nh)\geq0.
\]
Le coût correspond à l'oubli de la fibre, non au transport réversible de
l'état enrichi ; cette identité n'affirme ni une dissipation physique totale nulle ni
une complexité computationnelle réduite.

\subsection*{Clôture opératorielle : de l'unité nonadique à la dimension continue}

La seconde clôture réalise opératoriellement la conservation évolutive de
l'unité. Pour le système complet, on considère
\[
 A_m=M_{9^m}(\mathbb C),
 \qquad \iota_m(a)=a\otimes I_9,
 \qquad
 \tau_m(a)=9^{-m}\operatorname{Tr}(a).
\]
L'identité \(\tau_{m+1}\circ\iota_m=\tau_m\) conserve simultanément
l'unité et la trace. La limite inductive est l'UHF de type surnaturel
\(9^\infty=3^\infty\), et la clôture faible de sa représentation GNS traciale est
le facteur hyperfini \(R_{\rm hyp}\) de type \(II_1\). Les valeurs
\(\operatorname{rank}(P)/9^m\) se densifient et la trace des projecteurs dans
\(R_{\rm hyp}\) réalise tout l'intervalle \([0,1]\) : une dimension continue
émerge d'une tour d'inclusions discrètes qui préservent l'unité.

Cette construction ne se confond pas avec la branche marquée
\(j_{m,r}(a)=a\otimes p_r\). Les isométries d'une fibre individuelle engendrent en
norme les compacts sur la limite hilbertienne et, en clôture faible, un facteur
de type \(I_\infty\). La sélection d'un secteur conserve la marque, mais non
l'unité de la tour ; c'est le relèvement simultané des neuf fibres qui
produit la branche UHF--\(II_1\). On n'identifie pas non plus sans preuve la réalisation
UHF à la réalisation par produit croisé de l'horloge profinie : une
conjugaison canonique exige de déclarer l'action, la mesure, la moyennabilité,
l'ergodicité et l'isotropie. La logique orthomodulaire des projecteurs, les contextes
booléens, les états normaux et la normalisation entropique
\[
 S_\tau(9^m\rho)=S_{\rm vN}(\rho)-m\log9
\]
sont dérivés au sein de cette construction opératorielle et maintiennent séparés
l'état matriciel, sa densité traciale et sa publication statistique.

La topologie, la dimension et la composition ne se confondent pas dans le facteur. La clôture
transversale est enregistrée au moyen de
\[
 \mathfrak C_{\rm HMT}
 =\bigl(C(\Omega_\infty),R_{\rm hyp},\mathcal G_{\rm TPK}\bigr).
\]
Le spectre commutatif conserve les points et les cylindres ; le facteur tracial conserve
la dimension et la statistique ; le groupoïde conserve la concaténation, l'inversion
admissible et la provenance des routes. Cette architecture triple est formulée après
la double projection et le corridor exceptionnel. Elle ne les fonde ni ne les
remplace : l'évaluation archimédienne et la chaîne
Paley--Witt--Golay--Leech--\(V^\natural\)--Monstre--Moonshine demeurent
les deux lectures coordonnées du même antécédent enrichi.

\subsection*{Clôture matricielle : profondeur, taille extérieure et dilatation}

La troisième clôture introduit un contrôle simultané de trois indices. Si
\(m\) mesure la profondeur généalogique, \(n\) la taille extérieure d'une grille
et \(s\) le niveau matriciel de ses entrées, la famille
\(\mathfrak M_{m,s}^{(n)}\) est bigraduée par \((m,s)\) lorsque \(n\) demeure
fixe et trigraduée par \((m,n,s)\) lorsque la taille extérieure varie aussi.
En particulier,
\[
 \mathcal K_{m,s}
 :=\operatorname{UCP}\!\bigl(C(X_m),M_s(\mathbb C)\bigr)
\]
possède deux opérations indépendantes : la restriction généalogique
\(\rho_m^s(\Phi)=\Phi\circ\jmath_m\) et la compression matricielle
\[
 \operatorname{mc}_{\boldsymbol V}(\Phi_1,\ldots,\Phi_k)(f)
 =\sum_rV_r^*\Phi_r(f)V_r,
 \qquad \sum_rV_r^*V_r=I_s.
\]
Les deux commutent exactement :
\[
 \rho_m^s\!\left(\operatorname{mc}_{\boldsymbol V}(\Phi_1,\ldots,\Phi_k)\right)
 =\operatorname{mc}_{\boldsymbol V}\!\left(
   \rho_m^{s_1}(\Phi_1),\ldots,\rho_m^{s_k}(\Phi_k)\right).
\]
L'identité est fonctorielle et distingue raffinement de l'histoire et réduction
de l'observateur.

Toute famille compatible en profondeur détermine une unique application UCP
sur \(C(X_\infty)\) et admet une dilatation de Stinespring commune. Dans le
corridor diagonal, une unique mesure spectrale projective \(P\) satisfait
\[
 E_x^{(m)}=V^*P([x]_m)V
\]
pour toutes les profondeurs et tous les cylindres. La dilatation minimale est
unique à équivalence unitaire près et conserve exactement ce qui est vu par le
lecteur ; une somme directe permet une réalisation environnementale fidèle qui conserve
toutes les distinctions de l'algèbre. Sur la limite UHF, une famille compatible
d'applications UCP possède également un prolongement et une dilatation communs.
L'espérance conditionnelle
\[
 \mathbb E_m=\operatorname{id}\otimes\tau_9:
 M_{9^{m+1}}(\mathbb C)\longrightarrow M_{9^m}(\mathbb C)
\]
préserve la trace et écarte le dernier facteur de manière contrôlée ; son environnement
de Stinespring enregistre le chiffre omis. En général, il n'existe pas de mémoire
auxiliaire finie uniforme pour toutes les profondeurs. De plus, la fidélité d'une
représentation n'implique pas à elle seule la fidélité dynamique APP--TPK : celle-ci
exige un opérateur d'entrelacement compatible avec la route, la monodromie et la composition.

La théorie matricielle se concrétise en objets vérifiables. Les carrés magiques
quantiques qutrit de paramètres \((9,3)\) et \((12,3)\) possèdent respectivement
81 et 144 effets, des sommes de lignes et de colonnes exactement unitaires et des
dilatations simultanées déclarées ; les mêmes POVM produisent des cubes magiques
quantiques d'ordres neuf et douze. Le cas d'ordre douze est organisé au moyen de
la fibration \(12\to4\times3\), avec sa jauge et sa condition explicite
d'équivariance. Par une autre voie, les espérances conditionnelles somme et produit de l'APP
engendrent l'algèbre
\[
 C^*(P_{\Sigma,2},P_{\Pi,2})
 \cong\mathbb C^4\oplus M_2(\mathbb C)\oplus M_2(\mathbb C),
\]
et le bloc central réalise à chaque niveau \(s\) l'intervalle matriciel libre
\([5I_s,9I_s]\). Ces constructions distinguent carrés latins quantiques,
carrés magiques, mesures projectives, mesures positives et semi-classicité ;
aucune identification n'est obtenue uniquement parce que les objets partagent l'ordre
neuf.

L'antécédent spectral commun de deux lecteurs n'est formulé qu'après
les avoir construits tous les deux. Pour des applications finies
\(R_m:X_m\to\mathcal A_m\) et \(Q_m:X_m\to\mathcal E_m\), les incidences
\[
 P_{a,e}=1_{\{R_m=a\}}1_{\{Q_m=e\}}
\]
raffinent les deux mesures spectrales diagonales et récupèrent leurs marginales. Le
résultat ne prouve pas l'injectivité conjointe et n'impose pas qu'un lecteur soit
une compression de l'autre. Dans le régime non commutatif, deux PVM qutrit mutuellement
non biaisées ne possèdent pas de PVM conjointe nette ; la réalisation simultanée exige
un bruit compatible ou un environnement plus grand. Cette différence de type préserve la
double projection HMT : le théorème matriciel classe des réalisations de lecteurs
et ne réécrit pas leur antécédent fondateur.

\subsection*{Clôture physique : domaine commun, deux orientations et localisation}

La quatrième clôture construit le domaine depuis lequel une route TPK peut recevoir
simultanément une réalisation spectrale, hilbertienne, dimensionnelle et de Cartan.
Soient \(\mathsf D_{\rm sp}\), \(\mathsf D_Q\), \(\mathsf D_D\) et
\(\mathsf D_C\) les domaines de référence de ces quatre représentations, avec des
foncteurs d'oubli fidèles vers la catégorie
\(\mathsf P_{\rm TPK}^{\rm enr}\) des états préparés et des routes enregistrées.
Le produit fibré
\[
 \mathsf P_{\rm CT}^{+}
 :=\mathsf D_{\rm sp}
 \times_{\mathsf P_{\rm TPK}^{\rm enr}}\mathsf D_Q
 \times_{\mathsf P_{\rm TPK}^{\rm enr}}\mathsf D_D
 \times_{\mathsf P_{\rm TPK}^{\rm enr}}\mathsf D_C
\]
impose que les quatre composantes aient exactement le même antécédent.
Ses projections canoniques sont des adaptateurs fonctoriels et produisent
\[
 \mathfrak R_{\rm MD}^{\rm disc,+}
 =\bigl(\mathfrak S_+,\mathcal Q_+,\mathfrak D_+,
        \rho_{C,+}^{\rm disc}\bigr),
\]
à valeurs, respectivement, dans les routes spectrales, les espaces de Hilbert réels
à structure complexe, le monoïde dimensionnel positif et le groupe de
Poincaré discret. La construction affirme le foncteur sur le domaine commun ; elle
n'affirme pas que toute route TPK admette automatiquement les quatre relèvements.

Sur une route \(\gamma=e_1\cdots e_r\), la variation totale
\[
 V_+(\gamma)=\sum_jw(e_j)
\]
et le transport net orienté
\[
 V_{\rm or}(\gamma)=\sum_j\varepsilon(e_j)w(\underline e_j)
\]
sont des lecteurs frères. La boucle \(ee^\vee\) démontre que le premier ne se
factorise pas, en général, par le second : elle possède un transport orienté nul et une
variation positive. Pour obtenir une réalisation inversible, on restreint le
domaine aux arêtes dont les composantes spectrale et hilbertienne sont
inversibles, dont la donnée dimensionnelle est une cochaîne additive et dont les cocycles de
Cartan sont normalisés et antisymétriques. Ce n'est qu'alors que la localisation
\(\Pi_\vee(\mathsf P_{{\rm CT},{\rm rev}}^+)\) produit le foncteur orienté
\(\mathfrak R_{\rm MD}^{\rm disc,or}\) et transporte l'inversion composante par
composante. La lecture positive n'est pas éliminée et la lecture orientée n'est pas obtenue
en changeant uniquement un signe.

Cette construction physique précède les constructions de Planck et de l'électron
parce qu'elle leur fournit domaine, représentations et règles de transport, mais
n'absorbe pas leurs démonstrations de référence ni la loi générale des masses. La
chaîne nombre--particule--route--registre chronologique enrichi--signature--
caractère--tours conserve sa direction causale ; l'échelle absolue d'action,
l'électron central, les familles de masse, CKM, PMNS et les secteurs physiques
maintiennent leurs objets, domaines et conclusions propres. Une coïncidence métrologique n'est
confrontée qu'après avoir figé la sortie interne. Les extensions qui
exigent encore un opérateur d'entrelacement, une limite lisse, une preuve de stabilité dynamique
ou un morphisme physique sont énoncées avec ces données comme hypothèses exactes.
L'absence de l'une de ces applications empêche d'étendre le théorème au codomaine
correspondant, sans altérer les théorèmes généalogiques, opératoriels et
matriciels déjà démontrés dans leurs domaines.

La lecture de l'ouvrage suit, par conséquent, l'ordre de dépendance et non
celui de la reconnaissance extérieure :
\[
\begin{aligned}
 \mathrm{APP}&\longrightarrow\mathrm{TRIT}_{\rm orientado}
 \longrightarrow\mathrm{TPK}\\
 &\longrightarrow\text{section enrichie}
 \longrightarrow
 \begin{cases}
  \text{évaluation archimédienne},\\
  \text{projection exceptionnelle},\\
  \text{algèbres cylindriques et clôture traciale},\\
  \text{routes préparées et réalisation physique}.
 \end{cases}
\end{aligned}
\]
Chaque sortie conserve son domaine, ses hypothèses et ses obstructions précises. Ce qui est commun n'est pas
une égalité forcée entre codomaines, mais la provenance du même
état enrichi. Tel est le sens précis de la clôture holographique : aucune
publication terminale n'épuise la généalogie qui la produit et aucune clôture
ultérieure ne peut remplacer la double direction --archimédienne et exceptionnelle--
qui conserve l'unité avec mémoire.
